\section{Pohon Chow-Liu order-1 via Algoritma Prim (MemoriO(D))}

Struktur dependensi disusun sebagai pohon rentang maksimum menggunakan algoritma Prim, yang hanya menyimpan larik sepanjang $D$ (bukan matriks mutual information $D \times D$); mutual information dihitung baris-demi-baris atas representasi terkemas-bit (\texttt{modul codaeda/algos.py}).

\begin{verbatim}
def chow_liu_prim(A, rng, max_rows=1536, eps=1e-6):
    n, D = A.shape
    idx = rng.choice(n, max_rows, replace=False) 
    if n > max_rows else np.arange(n)

    # Bit-packing: kolom biner dikemas 8 baris/byte; 
    # ko-okurensi via AND bitwise + tabel popcount, 
    # tanpa membentuk matriks MI D x D maupun salinan 
    # float data.

    # (disusun per blok baris)
    P    = np.packbits(A[idx], axis=0) 
    cnt1 = _POPCNT[P].sum(axis=0).astype(float); 
    p1 = cnt1 / n

    # mutual information baris-ke-i terhadap 
    # semua variabel
    def mi_row(i):            
        j11 = _POPCNT[P & P[:, i:i+1]].sum(axis=0) / n
        # ... hitung j10, j01, j00 
        # lalu MI = SUM p*log2(p / (p_i * p_j)) ...
        return mi

    # Algoritma Prim: HANYA larik sepanjang D 
    # (bukan matriks D x D) -> memori O(D)
    parent = np.full(D, -1); best = np.full(D, -np.inf); 
    intree = np.zeros(D, bool)
    cur = 0; intree[cur] = True

    for _ in range(D - 1):
        # satu baris MI dihitung saat perlu
        m   = mi_row(cur)                       
        upd = (~intree) & (m > best); 
        best[upd] = m[upd]; bestfrom[upd] = cur
        nxt = int(np.argmax(np.where(intree, 
        -np.inf, best)))
        parent[nxt] = bestfrom[nxt]; intree[nxt] = True; 
        cur = nxt

    return parent, cp0, cp1, wgt   
    # induk + tabel probabilitas bersyarat per variabel
\end{verbatim}